\subsection{An explicit integer for a character orbit}
\label{selmer-resultant}

The elementary part of the congruence estimate can be expressed by a
single resultant. This separates the construction of a useful integer
from the arithmetic assertion that its valuations control Selmer lengths.

\begin{proposition}[A simultaneous integer certificate]
\label{selmer-resultant-certificate}
Let \(Q\ge5\), \(e\ge1\), and \(a=Q^e\). Let
\(\alpha_1,\ldots,\alpha_R\) be a Galois-stable multiset of
algebraic numbers with \(|\alpha_j|\le3\) and \(a\alpha_j\)
integral. Let \(z_1,\ldots,z_m\) be a Galois-stable multiset of
roots of unity. Put
\[
 P(T)=\prod_j(T-a\alpha_j),\qquad Z(X)=\prod_i(X-z_i),
 \qquad V(X)=X^R P\!\left(X+a+\frac{a^2}{X}\right).
\]
Then \(P,Z,V\) have integer coefficients and the integer
\(n_Q=\operatorname{Res}(Z,V)\) is nonzero. It satisfies
\[
 \log|n_Q|\le mR\log(a^2+4a+1)
                 \le mR(2e\log Q+\log2).
\]
For every prime \(\ell\nmid Q\), every embedding into
\(\overline{\mathbf Q}_\ell\), and all subsets
\(J_i\subseteq\{1,\ldots,R\}\),
\[
 \sum_i\sum_{j\in J_i}
 v_\ell\bigl(\alpha_j-(Q^{-e}z_i+1+Q^e z_i^{-1})\bigr)
                                      \le v_\ell(n_Q),
 \qquad v_\ell(\ell)=1.
\]
Multiplicity is retained throughout; empty products equal one.
\end{proposition}

\begin{proof}
Galois invariance and integrality give \(P,Z\in\mathbf Z[X]\).
Writing \(P(T)=\sum_{k=0}^R c_kT^k\), we have
\[
 V(X)=\sum_{k=0}^R c_kX^{R-k}(X^2+aX+a^2)^k\in\mathbf Z[X].
\]
The resultant product formula is
\[
 n_Q=\left(\prod_i z_i\right)^R
             \prod_{i,j}(z_i+a+a^2z_i^{-1}-a\alpha_j).
\]
The first factor is \(\pm1\). Each remaining factor has modulus
between \(a^2-4a-1>0\) and \(a^2+4a+1<2a^2\).
This proves nonvanishing and the logarithmic estimate. At
\(\ell\nmid Q\), the factors are integral and multiplication by
\(a\) is a unit. The valuation of the product is therefore the sum
of all the nonnegative trace-difference valuations. Discarding those
outside the subsets \(J_i\) proves the last assertion.
\end{proof}

\subsection{Combine two bounds for the same Selmer group}
\label{selmer-common-divisor}

\begin{proposition}[The two exceptional primes]
\label{selmer-corrected-gcd}
Fix distinct primes \(r,s\) and nonzero integers \(n_r,n_s\).
For each prime in a specified finite set \(\mathcal T\), let
\(L_\ell\ge0\) satisfy
\[
 L_\ell\le2v_\ell(n_r)\quad(\ell\ne r),\qquad
 L_\ell\le2v_\ell(n_s)\quad(\ell\ne s).
\]
Set \(A=|n_r|\), \(B=|n_s|\), and
\[
 G_{r,s}=
 \gcd\!\left(A/r^{v_r(A)},B/s^{v_s(B)}\right)
             r^{v_r(B)}s^{v_s(A)}.
\]
Then \(L_\ell\le2v_\ell(G_{r,s})\) on \(\mathcal T\).
For any prescribed weights \(0\le w_\ell\le1\),
\[
 \sum_{\ell\in\mathcal T}w_\ell L_\ell\log\ell
 \le2\sum_{\ell\in\mathcal T}w_\ell
                      v_\ell(G_{r,s})\log\ell
 \le2\log G_{r,s}\le2\log A+2\log B.
\]
If \(r,s\notin\mathcal T\), the same first inequality holds
with \(\gcd(A,B)\) instead of \(G_{r,s}\).
\end{proposition}

\begin{proof}
Away from \(r,s\), both estimates concern the same \(L_\ell\),
so their minimum applies. The valuation of \(G_{r,s}\) is that
minimum. Its valuations at \(r,s\) are respectively
\(v_r(B),v_s(A)\), exactly the estimates not excluded there.
Thus \(G_{r,s}\mid AB\), and prime factorization gives the weighted
inequalities.
\end{proof}

In the arithmetic application, \(L_\ell\) is the sum of normalized
Selmer lengths over one fixed character orbit, with the same host and
twist for both auxiliary primes. The required arithmetic estimate still
has to bound that length by twice the sum of the trace-difference
valuations in Proposition~\ref{selmer-resultant-certificate}.
The spectral families must classify the very representations used in
this estimate, injectively up to isomorphism for each character.
These conditions cannot be replaced by rational vanishing or a bound
for unrelated representations. Nor does the greatest common divisor of
integers bounding different Selmer components bound their sum.

When both comparisons are available, the common-divisor bound preserves
every actual prime weight and can be strictly smaller than the product
bound. Its quantitative size for the required arithmetic families is
not asserted here.
